% Fiche Python 3 (version ingénieur) — recto-verso A4 paysage
% Source : pas de cours ; périmètre choisi : le langage (recto) et les scripts d'ingénieur
%   (fichiers, CSV/JSON, argparse, subprocess, re, numpy, matplotlib, bits, pièges) au verso.
% Forme une paire avec python_avance.tex (décorateurs, générateurs, typing, asyncio...) :
%   ici seulement l'essentiel de ces sujets (itertools, logging, breakpoint), le détail est là-bas.
% Compiler : make (ici ou à la racine), ou xelatex python_cheat.tex — une seule passe suffit
% Apparence : ../librairie (fiche.sty + theme-python.tex). Macros de contenu :
%   colonne[titre][couleur] (3 par page, séparées par \colbreak) · bloc{titre}{cat}
%   catégories : syn typ ctl ex io lib (couleurs et icônes dans le thème)
%   code · \V|...| code en ligne · «...» dans le code = à remplir
%   \sub sous-titre · \note remarque · \piege piège · \badge{TXT}
%   optab étiquette | texte · tab3 étiquette | code | remarque
%   (largeurs mesurées par la librairie : celles écrites entre [] et {} sont ignorées)
% Le code est en police à chasse fixe : les commentaires sont alignés à la main
% (attention : les « » ne sont pas affichés, ils ne comptent pas dans l'alignement).
% Après compilation : grep "\[fiche\]" python_cheat.log (place libre, collisions).
\documentclass{article}
\usepackage[python]{../librairie/fiche}   % apparence : ../librairie (fiche.sty + theme-python.tex)
% --- Langage Python 3 pour listings ---
% Classes de mots-clés : [1] rose · [2] cyan · [3] vert · [4] couleur du thème
% (styles dans ../librairie/fiche.sty). Chargé par python_cheat.tex.
% [1] mots-clés · [2] types · [3] fonctions / modules intégrés
\lstdefinelanguage{CODE}{
  sensitive=true, alsoletter={@},
  morekeywords=[1]{False,None,True,and,as,assert,async,await,break,class,continue,def,
    del,elif,else,except,finally,for,from,global,if,import,in,is,lambda,nonlocal,not,
    or,pass,raise,return,try,while,with,yield,match,case},
  morekeywords=[2]{int,float,complex,bool,str,bytes,list,tuple,dict,set,frozenset,
    range,object,Path,ndarray,Exception,ValueError,KeyError,IndexError,TypeError,
    FileNotFoundError,ZeroDivisionError,OSError,RuntimeError,StopIteration,
    defaultdict,Counter,deque},
  morekeywords=[3]{print,len,enumerate,zip,map,filter,sorted,reversed,sum,min,max,abs,
    round,open,input,isinstance,type,super,any,all,iter,next,repr,id,hash,getattr,
    hasattr,setattr,divmod,bin,hex,ord,chr,self,cls,@dataclass,@property,
    __init__,__repr__,__name__,__doc__,
    np,plt,os,sys,re,json,csv,math,subprocess,argparse,shutil,itertools,pathlib,
    random,time,logging,statistics,struct,datetime},
  morecomment=[l]{\#}, morestring=[b]", morestring=[b]',
}
   % coloration syntaxique : \lstdefinelanguage{CODE}

\begin{document}
%%%%%%%%%%%%%%%%%%%%%%%%%%%%%%%%%%%%%%%%%%%%%%%%%%%%%%%%%%%%%%%%%%%%%%%%%%%%%%
%                         RECTO : LE LANGAGE
%%%%%%%%%%%%%%%%%%%%%%%%%%%%%%%%%%%%%%%%%%%%%%%%%%%%%%%%%%%%%%%%%%%%%%%%%%%%%%
\entete{Python 3}{recto · le langage}
\noindent
% ======================= colonne 1 : bases =======================
\begin{colonne}[bases][col-a]
\begin{bloc}{Opérateurs}{syn}
\begin{tab3}[13mm]{26mm}
Arithm.   & \V|+ - * / // % **| & \V|/| réel, \V|//| entier, \V|**| puissance \\
Compar.   & \V|== != < <= > >=|   & chaînables : \V|0 <= i < n| \\
Logiques  & \V|and or not|        & court-circuit, rendent un opérande \\
Identité  & \V|is|, \V|is not|    & même objet (pour \V|None|) \\
Appart.   & \V|in|, \V|not in|    & liste, chaîne, clé de dict \\
Bits      & \V/& | ^ ~ << >>/     & sur les \V|int| \\
Ternaire  & \V|a if «c» else b|   & une expression \\
Morse     & \V|(n := len(x))|     & affecte dans une expression \badge{3.8}
\end{tab3}
\end{bloc}

\begin{bloc}{Types et nombres}{typ}
\begin{code}
n = 42           # int, sans limite de taille
x = 3.14         # float : double IEEE 754
z = 2 + 3j       # complex
ok = True        # bool (sous-classe d'int)
s = "texte"      # str : Unicode, immuable
b = b"\x01\xff"  # bytes
r = None         # absence de valeur
type(x), isinstance(x, (int, float))
\end{code}
\note{Immuables : \V|int float str tuple bytes frozenset|. Mutables : \V|list dict set| et la plupart des objets.}
\sub{Nombres}
\begin{code}
int("ff", 16), int("0b101", 0)  # 255, 5
hex(255), bin(5), oct(8)        # '0xff', '0b101', '0o10'
-7 // 2, -7 % 2, divmod(17, 5)  # -4, 1, (3, 2)
round(2.5), round(3.14159, 2)   # 2 (au pair), 3.14
int(3.9), abs(-2), 2 ** 0.5     # 3 (tronque), 2, 1.414
1_000_000, 1e-9, float("inf")
import math                     # pi, sqrt, log2, ceil...
\end{code}
\end{bloc}

\begin{bloc}{Chaînes et f-strings}{typ}
\begin{schema}[x=4.2mm, y=4.2mm, font=\fontcode]
  \foreach \c [count=\i from 0] in {P,y,t,h,o,n}
    \node[draw=trait, line width=0.5pt, fill=cat!\ifnum\i>0 \ifnum\i<4 35\else 8\fi\else 8\fi,
          minimum size=4.2mm, inner sep=0pt] at (\i+0.5,0) {\c};
  \foreach \i in {0,...,6} \node[text=cat!80!black, inner sep=0.5pt] at (\i,0.95) {\i};
  \foreach \i in {0,...,5} \node[text=mk-grey, inner sep=0.5pt] at (\i,-0.95) {\the\numexpr\i-6\relax};
  \node[anchor=west, text=note, font=\fontbloc, align=left] at (6.6,0)
    {\V|s[1:4]| $\to$ \V|'yth'| (fin exclue)\\ \V|s[-2:]| $\to$ \V|'on'|};
\end{schema}
\begin{code}
s[0], s[-1], s[1:4], s[::-1]      # [::-1] : inverse
"a,b,c".split(",")                # ['a', 'b', 'c']
", ".join(["a", "b"])             # 'a, b'
s.strip(), s.lower(), s.upper()
s.replace("a", "b"), s.find("x")  # -1 si absent
s.startswith("x"), "ab" in s
f"{x:.3f} {n:5d} {n:#06x}"        # '3.140    42 0x002a'
f"{s:>10} {s:<10} {s:^10}"        # aligner sur 10
f"{x=}"                           # 'x=3.14' (débogage)
r"C:\new"                         # brute : \ non échappé
\end{code}
\end{bloc}

\begin{bloc}{Console}{io}
\begin{code}
print(a, b, sep=", ", end="")  # sans retour ligne
print(f"{nom:<8}{v:6.2f}")     # colonnes alignées
nom = input("nom ? ")          # toujours une str
n = int(input("n ? "))
import sys; sys.argv, sys.exit(1)  # arguments, code
\end{code}
\note{Fichiers et arguments nommés : verso (colonne 4).}
\end{bloc}
\end{colonne}\colbreak
% ======================= colonne 2 : contrôle, fonctions, modules =======================
\begin{colonne}[contrôle · fonctions][col-b]
\begin{bloc}{Contrôle et boucles}{ctl}
\begin{code}
if x > 0:
    ...
elif x == 0:
    ...
else:
    ...
while «cond»:
    if «fini»: break  # sort de la boucle
    continue        # tour suivant
for i in range(«début», «fin», «pas»): ...
for i, x in enumerate(l, start=1): ...
for a, b in zip(l1, l2): ...  # s'arrête au plus court
for x in reversed(l): ...
any(x > 0 for x in l), all(l)
min(l, key=len), max(d, key=d.get)
match cmd.split():  # filtrage de motifs (3.10)
    case ["go", dest]: aller(dest)
    case ["stop"] | ["q"]: quitter()
    case _: print("inconnu")
\end{code}
\note{Les blocs sont définis par l'indentation (4 espaces). \V|pass| = bloc vide.}
\end{bloc}

\begin{bloc}{Fonctions}{ctl}
\begin{code}
def f(a, b=2, *args, k=0, **kw) -> int:
    """Docstring : ce que fait f."""
    return a + b               # None si pas de return
f(1); f(1, 3, 4, 5, k=1, x=9)  # args=(4, 5), kw={'x': 9}
f(*liste, **dico)              # dépaqueter à l'appel
carre = lambda x: x * x        # fonction anonyme
def g(x: float, y: list[int] | None = None): ...
\end{code}
\note{Paramètres après \V|*args| (ou \V|*|) : seulement par mot-clé. Les annotations ne sont pas vérifiées à l'exécution.
Affecter une variable dans une fonction la rend locale : \V|global| / \V|nonlocal| pour modifier celle du dessus.}
\piege{\V|def f(l=[])| garde la même liste d'un appel à l'autre : défaut \V|None|, puis \V|l = l or []|.}
\end{bloc}

\begin{bloc}{Exceptions}{ctl}
\begin{code}
try:
    v = int(txt)
except (ValueError, TypeError) as e:
    print("invalide :", e)
else:
    ...  # si aucune exception
finally:
    ...  # toujours exécuté
raise ValueError(f"mauvais n : {n}")
class SimError(Exception): pass
\end{code}
\note{\V|with| ferme le fichier même en cas d'exception.}
\end{bloc}

\begin{bloc}{Modules et imports}{lib}
\begin{code}
import numpy as np        # alias
from pathlib import Path  # un nom
from . import utils       # relatif (paquet)
if __name__ == "__main__": main()
\end{code}
\note{Fichier = module ; dossier avec \V|__init__.py| = paquet.}
\end{bloc}
\end{colonne}\colbreak
% ======================= colonne 3 : collections, classes =======================
\begin{colonne}[collections · classes][col-c]
\begin{bloc}{Listes et tuples}{lib}
\begin{code}
l = [3, 1, 2]
l.append(4); l.extend([5, 6])  # en place
l.insert(0, 9); l.pop(); l.pop(0)
l.remove(3)                    # 1re occurrence
l.sort(key=abs, reverse=True)  # en place, rend None
sorted(l)                      # nouvelle liste triée
l.index(2), l.count(2), len(l)
l[1:3] = []                    # suppr. une tranche
t = (1, "a")                   # tuple : immuable
a, *reste = [1, 2, 3]          # a=1, reste=[2, 3]
a, b = b, a                    # échange
\end{code}
\piege{\V|b = a| ne copie pas : \V|a.copy()|, ou \V|copy.deepcopy(a)| si imbriqué.}
\end{bloc}

\begin{bloc}{Dictionnaires et ensembles}{lib}
\begin{code}
d = {"clk": 10, "rst": 0}
d["en"] = 1                   # ajout / mise à jour
d.get("x", 0)                 # défaut si absente
d.setdefault("k", []).append(1)
for k, v in d.items(): ...    # .keys() .values()
d.pop("rst"); d | {"a": 1}    # retirer, fusionner
"clk" in d                    # test de clé
e = {1, 2, 3}; e.add(4)
e | f, e & f, e - f, e ^ f    # union, inter, diff
from collections import defaultdict, Counter
cnt = Counter("abracadabra")  # {'a': 5, 'b': 2...}
grp = defaultdict(list)       # grp[k].append(v)
\end{code}
\note{Clés : objets hachables (immuables). L'ordre d'insertion est conservé.}
\end{bloc}

\begin{bloc}{Compréhensions}{lib}
\begin{code}
[x**2 for x in range(10) if x % 2]  # liste
[x if x > 0 else 0 for x in l]      # avec ternaire
{k: v for k, v in zip(noms, vals)}  # dict
{c.lower() for c in s}              # ensemble
sum(x**2 for x in data)             # générateur
[[0] * n for _ in range(m)]         # matrice m x n
[x for ligne in m for x in ligne]   # aplatir
\end{code}
\end{bloc}

\begin{bloc}{Classes}{ctl}
\begin{code}
class Reg:
    nb = 0            # attribut de classe
    def __init__(self, nom, w=8):
        self.nom, self.w = nom, w
        Reg.nb += 1
    def __repr__(self):
        return f"Reg({self.nom!r}, {self.w})"
    @property
    def maxi(self): return 2**self.w - 1
class Compteur(Reg):  # héritage
    def __init__(self, nom):
        super().__init__(nom, 16)
\end{code}
\end{bloc}
\end{colonne}

%%%%%%%%%%%%%%%%%%%%%%%%%%%%%%%%%%%%%%%%%%%%%%%%%%%%%%%%%%%%%%%%%%%%%%%%%%%%%%
%                     VERSO : SCRIPTS D'INGÉNIEUR
%%%%%%%%%%%%%%%%%%%%%%%%%%%%%%%%%%%%%%%%%%%%%%%%%%%%%%%%%%%%%%%%%%%%%%%%%%%%%%
\newpage
\entete{Python 3}{verso · scripts d'ingénieur}
\noindent
% ======================= colonne 4 : entrées, sorties, fichiers =======================
\begin{colonne}[fichiers · scripts][col-a]
\begin{bloc}{Fichiers texte}{io}
\begin{code}
with open("in.txt", encoding="utf-8") as fh:
    for ligne in fh:           # ligne par ligne
        ligne = ligne.rstrip("\n")
with open("out.txt", "w") as fh:  # "a" = ajouter
    fh.write(f"{n}\n")
    print("ok", file=fh)
txt = Path(«f»).read_text()      # tout, fermé aussitôt
lignes = txt.splitlines()
data = Path(«f»).read_bytes()    # binaire
\end{code}
\note{Modes : \V|"r"| lecture, \V|"w"| écrase, \V|"a"| ajoute, \V|"b"| binaire (bytes).}
\end{bloc}

\begin{bloc}{Chemins (pathlib)}{io}
\begin{code}
from pathlib import Path
p = Path("runs") / "synth" / "area.rpt"
p.exists(), p.is_file(), p.is_dir()
p.name, p.stem, p.suffix  # area.rpt area .rpt
p.parent, p.with_suffix(".csv")
p.read_text(); p.write_text("...")
Path("out").mkdir(parents=True, exist_ok=True)
for f in Path("runs").glob("*/area.rpt"): ...
for f in Path(".").rglob("*.v"): ...  # récursif
Path.cwd(), Path.home(), p.resolve()
\end{code}
\end{bloc}

\begin{bloc}{CSV et JSON}{io}
\begin{code}
import csv, json
with open("res.csv", newline="") as fh:
    for row in csv.DictReader(fh):  # dict par ligne
        f = float(row["freq"])      # valeurs en str
with open("out.csv", "w", newline="") as fh:
    w = csv.writer(fh)
    w.writerow(["design", "area"])
    w.writerows(lignes)
cfg = json.loads(Path("cfg.json").read_text())  # dict
Path("o.json").write_text(json.dumps(cfg, indent=2))
import pandas as pd
df = pd.read_csv("res.csv"); df["area"].mean()
\end{code}
\end{bloc}

\begin{bloc}{Script : argparse, main}{ex}
\begin{code}
#!/usr/bin/env python3
"""Rassemble les rapports de synthèse."""
import argparse, sys
def main() -> int:
    ap = argparse.ArgumentParser(description=__doc__)
    ap.add_argument("dossier")  # positionnel
    ap.add_argument("-o", "--out", default="res.csv")
    ap.add_argument("-n", type=int, default=10)
    ap.add_argument("-v", action="store_true")
    args = ap.parse_args()      # args.out...
    return 0
if __name__ == "__main__":      # pas à l'import
    sys.exit(main())
\end{code}
\note{\V|python3 script.py runs -o a.csv -v| ; \V|-h| est généré.}
\end{bloc}
\end{colonne}\colbreak
% ======================= colonne 5 : outils =======================
\begin{colonne}[outils][col-b]
\begin{bloc}{Lancer des outils (subprocess)}{lib}
\begin{code}
import subprocess, os, shutil
r = subprocess.run(["vsim", "-c", "-do", "run.do"],
                   capture_output=True, text=True,
                   cwd="sim", timeout=600)
r.returncode, r.stdout, r.stderr
subprocess.run(«cmd», check=True)  # exception si échec
os.environ.get("HOME"); os.environ["SEED"] = "3"
shutil.copy(«src», «dst»); shutil.rmtree(«dossier»)
shutil.which("vivado")           # chemin ou None
\end{code}
\note{Passer une liste d'arguments plutôt qu'une chaîne avec \V|shell=True| : pas de souci d'espaces ni d'injection.}
\end{bloc}

\begin{bloc}{Expressions régulières (re)}{lib}
\begin{code}
import re
m = re.search(r"Total area:\s+([\d.]+)", txt)
if m: area = float(m.group(1))
re.findall(r"\d+", "a1 b22 c333")  # ['1','22','333']
re.sub(r"\s+", " ", txt)           # remplacer
pat = re.compile(r"(?P<sig>\w+)\s*=\s*(?P<v>\d+)")
for m in pat.finditer(txt):
    print(m["sig"], int(m["v"]))
\end{code}
\begin{optab}[17mm]
\V|. \d \w \s| & tout caractère, chiffre, mot, blanc \\
\V|* + ? {n,m}| & 0+, 1+, 0 ou 1, de n à m (\V|*?| non gourmand) \\
\V|^ $ \b| & début, fin, limite de mot \\
\V|[a-f] (...)| & classe de caractères, groupe capturant
\end{optab}
\note{Toujours des chaînes brutes \V|r"..."| pour les motifs.}
\end{bloc}

\begin{bloc}{Exemple : rapports de synthèse en CSV}{ex}
\begin{code}
import csv, re
from pathlib import Path
RE_AREA = re.compile(r"Total cell area:\s+([\d.]+)")
RE_PWR = re.compile(r"Total Dynamic Power\s+=\s+([\d.]+)")
rows = []
for rpt in sorted(Path("runs").glob("*/report.txt")):
    txt = rpt.read_text()
    a, p = RE_AREA.search(txt), RE_PWR.search(txt)
    if not (a and p):
        print(f"incomplet : {rpt}"); continue
    rows.append([rpt.parent.name, float(a[1]),
                 float(p[1])])
with open("bilan.csv", "w", newline="") as fh:
    w = csv.writer(fh)
    w.writerow(["config", "area_um2", "power_mW"])
    w.writerows(rows)
\end{code}
\end{bloc}

\begin{bloc}{Modules utiles}{lib}
\begin{tab3}[16mm]{31mm}
math       & \V|sqrt log2 ceil pi| & \V|isclose(a, b)| \\
statistics & \V|mean median stdev|       & sur une liste \\
itertools  & \V|product combinations|    & \V|chain groupby| \\
random     & \V|seed randint choice|     & \V|random()| dans $[0,1[$ \\
time       & \V|perf_counter sleep|      & chronométrer \\
datetime   & \V|datetime.now()|          & \V|.strftime("%H:%M")| \\
logging    & \V|basicConfig info| & plutôt que \V|print|
\end{tab3}
\end{bloc}
\end{colonne}\colbreak
% ======================= colonne 6 : numpy, matplotlib, pièges =======================
\begin{colonne}[numpy · tracés · pièges][col-c]
\begin{bloc}{numpy}{lib}
\begin{code}
import numpy as np
a = np.array([1, 2, 3], dtype=float)
np.zeros((3, 4)), np.ones(5), np.eye(3)
np.arange(0, 1, 0.1), np.linspace(0, 1, 11)
a.shape, a.dtype, a.reshape(3, 1), a.T
a * 2 + 1, a @ b, np.sqrt(a)  # élément par élément
m[1, :], m[:, 0], m[m > 0]    # tranche, masque
a.mean(), a.std(), a.max(), a.argmax()
m.sum(axis=0)                 # somme par colonne
d = np.loadtxt("w.csv", delimiter=",", skiprows=1)
np.savetxt("o.csv", d, delimiter=",", fmt="%.4g")
\end{code}
\note{Les tranches sont des vues (pas des copies) : \V|a.copy()| pour copier. Broadcasting : \V|(3,1) + (4,)| donne \V|(3,4)|.}
\end{bloc}

\begin{bloc}{matplotlib}{lib}
\begin{code}
import matplotlib.pyplot as plt
fig, ax = plt.subplots(figsize=(6, 3.5))
ax.plot(t, v, label="v(t)")  # courbe
ax.semilogy(f, err, "o--")   # log en y
ax.set_xlabel("t (ns)"); ax.set_ylabel("V")
ax.grid(True); ax.legend()
fig.tight_layout(); fig.savefig("v.pdf")
plt.show()                   # fenêtre
\end{code}
\end{bloc}

\begin{bloc}{Bits et registres}{ex}
\begin{code}
(x >> 4) & 0xF               # extraire bits 7..4
x | (1 << 3), x & ~(1 << 3)  # mettre / effacer bit 3
f"{x:08b}", f"{x:04X}"       # binaire / hexa sur n
def signe(x, w):             # complément à 2
    return x - (1 << w) if x >> (w - 1) else x
x.bit_length(), bin(x).count("1")
x.to_bytes(4, "little"), int.from_bytes(b, "big")
import struct; struct.pack("<If", 7, 1.5)
\end{code}
\end{bloc}

\begin{bloc}{Environnement}{io}
\begin{code}
python3 -m venv .venv      # env. isolé
source .venv/bin/activate  # (Windows : Scripts)
pip install numpy matplotlib
pip freeze > requirements.txt
python3 -i script.py       # console après
breakpoint()               # débogueur pdb
\end{code}
\end{bloc}

\begin{bloc}{Pièges}{ex}
\begin{optab}[19mm]
Liste de listes & \V|[[0]*3]*3| répète la même ligne : passer par une compréhension. \\
Flottants & \V|0.1 + 0.2 != 0.3| : comparer avec \V|math.isclose|. \\
\V|is| ou \V|==| & \V|is| teste l'identité : seulement pour \V|None|. \\
Itérer et modifier & retirer d'une liste pendant un \V|for| saute des éléments : itérer sur une copie.
\end{optab}
\note{Défaut mutable : bloc Fonctions · alias : bloc Listes (recto).}
\end{bloc}
\end{colonne}
\end{document}
